\subsection*{Three new mechanisms on the original physical source}
The present continuation adds Sections~\ref{sec:nm-kernel},
\ref{sec:bf-actual-flux}--\ref{sec:zf-zero-offset} and
\ref{sec:xd-cross-curve}. The complete preceding weak-roof and singular
multiplier theorems in Sections~\ref{sec:fr-full-roof} and
\ref{sec:sm-singular-multiplier} are inherited.

Fix one actual flight $j$ and an unvisited marked disk.
Its strict perpendicular clearance $u$, direction $\vartheta$
and side $\sigma$ give the original physical angular source.
Theorem~\ref{thm:nm-kernel} proves a uniform bound $H_{\rm nm}$
for its momentum kernel at every forward or reversed collision.
Both the first-hit and final true-collision derivatives cancel the
respective incidence loss; exactly $l-2$ intermediate genuine collisions
use the proved native curve-growth bound. For the original joint
record selector $\Pi$ and every additional selector in $[0,1]$,
Corollary~\ref{cor:nm-joint} gives
\[
 \int h\one_\Pi N_E(T^j x)\one_{\{p_i\in I\}}\,d\nu_R^*(x)
       \le {H_{\rm nm}|E||I|\over4\pi R c_*}.
\]
Both source signs, all regular roots, the original section and the
single conditional normalization are already included.
A physical roof interval can be selected by taking $h$ supported
in that interval; the estimate does not assert roof essential height.
The unchanged insertion
\[
 I_m=1+\sum_{i=1}^{m-1}{q_0^{\min(i,m-i)}\over c_i},\qquad q_0=47/53,
\]
has marked moments for every $1\le\gamma<2$
(Corollary~\ref{cor:nm-moments}). At $u=0$,
Theorem~\ref{thm:nm-trace} constructs the positive-side strict physical
trace. It takes the actual one-sided limits of the original finite
section and roof-interval tests. The marked disk remains an ordinary
possible later scatterer; no collision is concealed.

The second mechanism computes the scalar boundary charge.
Writing $g=\nabla F$, $G=D^2F$ and $\theta$ for the Gauss angle,
Theorem~\ref{thm:bf-flux} identifies its density with
\[
 f_\varepsilon\,|d\theta|
 = {W_\varepsilon\over4\pi R c_*}
                  \left|\partial_t u\big|_\theta\right|d\vartheta .
\]
The physical roof vector is $G^{-1}g/(g^TG^{-1}g)$.
At the marked foot, the two true optical products give
\[
 \left|\det {\,\partial(p_0,p_m)\over\partial(\vartheta,u)}\right|
     = A_-^p A_+^p(h_-+h_+).
\]
The signed Wronskian in Corollary~\ref{cor:bf-wronskian} is retained
before taking absolute values. Original section cuts carry their
own $(\alpha,p)$ Jacobian.

Theorem~\ref{thm:zf-flux} addresses a genuine critical boundary
mechanism. Fix the specified clearance level $u_0\in[0,3/200]$ and a
closed physical critical word $q_*$. Its smooth marked branch $B$
is measured from that level. On the covered zero-offset components,
$B(q_*)=0$ and the actual selected contact belongs to the common
inverse-gradient chart of Lemma~\ref{lem:zf-formal}.
For $0<|g|<b_{\rm an}$, the exact identity on the boundary $B(g)=0$ is
\[
 DB(g)g=\int_0^1sD^2B(sg)[g,g]\,ds .
\]
The common formal analytic extension gives
$\|D_g^2B\|\le K_{Bg}(47/50)^{d'_j}$, with both terms in the
second derivative of a composition. Its quadratic factor cancels
the physical roof-vector denominator. Consequently
\[
 \mathfrak B^{\rm crit}_{j,u_0}\le C_{\rm zf}(47/50)^{d'_j},
 \qquad
 \sum_{d'_j\ge D}\mathfrak B^{\rm crit}_{j,u_0}
       \le {2C_{\rm zf}(47/50)^D\over1-47/50}.
\]
The complete physical word sum and arbitrarily short or disconnected
covered pieces are included before this uniform estimate.
For every fixed $0<\theta_*<1/2$, the deep sum at
$D=\lceil\theta_*m\rceil$ vanishes after multiplication by $m^2$.
Formal intermediate points are not called physical orbits.
Nonzero offsets retain the term $-B(q_*)/(g^TG^{-1}g)$;
uncovered contacts, the critical fan, section cuts and paired
singular endpoints are not removed.

Third, Theorem~\ref{thm:xd-cross-distortion} compares two actual
trajectories with the same radius, lifted center word and specified
matching interior homogeneity strips. Each is compared to the same
global disk-product minimum by the retained localization theorem.
No connecting physical path is required. Summable contact
discrepancies, the logarithmic incidence bound on closed strips and
a contracting slope recursion prove
\[
 C_{\rm cross}^{-1}\le
 {c_m(Q)J_m^u(Q)\over c_m(\widehat Q)J_m^u(\widehat Q)}
       \le C_{\rm cross},
\]
uniformly in the word and collision count. Fixing the terminal strip
also gives unweighted Jacobian comparison. Reversal gives the exact
$j_m/c_0$ normalization. Corollary~\ref{cor:xd-normalized-cocycles}
controls the optical weight divided by $c_0c_m$ with factor
$30000^2C_{\rm cross}$, paid once at each compared point rather than
per collision. Lemma~\ref{lem:xd-physical-terminal-limit} supplies a
genuine terminal-grazing family in the triangular table and verifies
every intervening obstacle clearance. The endpoint factor is necessary.

These mechanisms remove specified losses in the native marked source
and transverse comparison. They do not establish the unrestricted
pressure sum, roofwise comparison-pool capacity or
\eqref{eq:sc-original-obligation}. The complete physical sums,
the supremum over $(R,n,k)$, roof essential supremum and original
outer $B$ limit remain in the target. The following inherited
architecture records the geometric, arithmetic, conditional and
analytic interfaces to which the new results apply.
